\subsection{Escolha da rota experimental}

A definição da rota experimental partiu da comparação entre diferentes possibilidades de síntese de sabão, avaliando-se as variáveis independentes de cada uma. Inicialmente foram propostas três rotas distintas, posteriormente confrontadas por meio de uma matriz de decisão que ponderou critérios técnicos, operacionais e ambientais. A partir desse resultado, definiu-se a rota a ser executada e desenvolveram-se os cálculos necessários para sua realização.

\subsubsection{Rotas propostas}

Foram consideradas três rotas de produção, combinando diferentes matérias-primas, processos e formas finais do produto:

\begin{itemize}
    \item \textbf{A:} produção de sabão em barra por saponificação a quente de óleo de coco;
    \item \textbf{B:} produção de sabão em barra por saponificação a frio de óleo de soja;
    \item \textbf{C:} produção de sabão líquido por saponificação a quente de óleo de rícino.
\end{itemize}

\subsubsection{Matriz de decisão}

Cada rota foi avaliada segundo seis critérios: custo de matéria-prima, simplicidade operacional, segurança, viabilidade laboratorial, rendimento e impacto ambiental. Dessa forma Atribuindo a cada um uma pontuação de 1 a 5. A Tabela~\ref{tab:matriz_decisao} apresenta as notas e suas justificativas.

\begin{table}[H]
    \centering
    \footnotesize
    \caption{Matriz de decisão das rotas de produção propostas.}
    \label{tab:matriz_decisao}
    \begin{tabularx}{\textwidth}{>{\raggedright\arraybackslash}p{2.3cm} >{\raggedright\arraybackslash}X >{\raggedright\arraybackslash}X >{\raggedright\arraybackslash}X}
        \toprule
        \textbf{Critério} & \textbf{Rota A} (coco, quente, barra) & \textbf{Rota B} (soja, frio, barra) & \textbf{Rota C} (rícino, quente, líquido) \\
        \midrule
        Custo de matéria-prima & \textbf{3}: óleo de coco tem custo moderado & \textbf{5}: óleo de soja costuma ser mais barato e amplamente disponível & \textbf{2}: óleo de rícino geralmente é mais caro, elevando o custo final \\
        Simplicidade operacional & \textbf{3}: processo a quente exige aquecimento e maior controle térmico & \textbf{4}: processo a frio é mais simples na bancada e demanda menos equipamento & \textbf{3}: exige aquecimento e controle de viscosidade/diluição no sabão líquido \\
        Segurança (toxicidade) & \textbf{3}: uso de NaOH e aquecimento aumenta o risco no manuseio & \textbf{4}: ainda usa NaOH, mas sem aquecimento intenso o risco operacional é menor & \textbf{3}: uso de KOH com aquecimento mantém risco químico \\
        Viabilidade laboratorial & \textbf{5}: execução direta e tempo curto para avaliação & \textbf{2}: viável, porém com tempo de cura longo (4 a 6 semanas), o que limita o cronograma experimental & \textbf{3}: viável, mas com alta complexidade operacional devido ao controle constante de viscosidade e ajuste de processo \\
        Rendimento & \textbf{4}: boa conversão e boa dureza de barra no processo a quente & \textbf{2}: conversão aceitável, mas sabão tende a ser mais macio e com cura mais longa & \textbf{4}: boa performance para sabão líquido, com boa solubilidade e espuma \\
        Impacto ambiental & \textbf{3}: maior consumo de energia devido ao aquecimento & \textbf{4}: menor gasto energético por ser processo a frio, com possibilidade de perda de produto por rancificação & \textbf{3}: também demanda energia térmica e não utiliza matéria-prima residual \\
        \midrule
        \textbf{Total} & \textbf{21} & \textbf{21} & \textbf{18} \\
        \bottomrule
    \end{tabularx}
\end{table}

Com base na matriz de decisão, as rotas A e B empataram com 21 pontos, enquanto a rota C obteve 18 pontos. Assim, a escolha técnica mais adequada foi combinar os pontos fortes das duas rotas mais bem avaliadas.

A rota definida para execução experimental foi a \textbf{saponificação a quente com mistura de óleo de coco e óleo de soja}, seria uma nova opção AB. Essa decisão considera que a rota B, apesar de bem pontuada em custo e simplicidade, apresenta menor viabilidade laboratorial devido ao tempo de cura mais longo do processo a frio. Ao optar pela rota a quente, o experimento torna-se mais compatível com o cronograma acadêmico.

Espera-se obter um sabão com dureza moderada, boa formação de espuma e bom potencial de limpeza, além de índice de saponificação relativamente alto, de modo a garantir conversão eficiente da matéria-prima.

\subsubsection{Índice de saponificação}

O índice de saponificação é definido como a quantidade, em miligramas, de hidróxido de potássio (KOH) necessária para saponificar um grama de uma determinada amostra de óleo ou gordura. Sua determinação experimental é realizada por titulação, empregando os seguintes reagentes:

\begin{itemize}
    \item óleo (amostra a ser medida);
    \item solução de hidróxido de potássio (KOH);
    \item fenolftaleína;
    \item ácido para neutralização do KOH, geralmente HCl.
\end{itemize}

\noindent além das seguintes vidrarias e equipamentos:

\begin{itemize}
    \item erlenmeyer;
    \item condensador de refluxo;
    \item proveta.
\end{itemize}

O procedimento de determinação segue as etapas abaixo:

\begin{enumerate}
    \item adicionar uma massa conhecida de óleo no erlenmeyer;
    \item adicionar solução de hidróxido de potássio em excesso no erlenmeyer;
    \item acoplar o erlenmeyer ao condensador de refluxo;
    \item manter a solução em ebulição branda por 30 minutos;
    \item após esse período, deixar a solução esfriar por cerca de 5 minutos;
    \item adicionar o ácido de neutralização à proveta;
    \item adicionar algumas gotas de fenolftaleína ao erlenmeyer (cerca de duas gotas);
    \item titular a solução do erlenmeyer até o desaparecimento da coloração rosa.
\end{enumerate}

Alternativamente, o índice de saponificação de uma mistura de óleos pode ser estimado a partir de uma média ponderada dos valores reportados na literatura para cada óleo. A Tabela~\ref{tab:is_literatura} reúne os valores compilados para o óleo de coco e o óleo de soja, bem como a média adotada para cada um.

\begin{table}[H]
    \centering
    \footnotesize
    \caption{Índices de saponificação do óleo de coco e do óleo de soja segundo a literatura.}
    \label{tab:is_literatura}
    \begin{tabularx}{\textwidth}{l >{\raggedright\arraybackslash}X c}
        \toprule
        \textbf{Tipo de óleo} & \textbf{Índices de saponificação da literatura (mg KOH/g)} & \textbf{Média (mg KOH/g)} \\
        \midrule
        Óleo de coco & 259; 263; 271; 270 \cite{ghani2018physicochemical}; 249 \cite{perera2020determination}; \newline 252,67 \cite{chen2024multifrequency}; 248,12; 251,23 \cite{jans2025optimization} & 255,84 \\
        Óleo de soja & 195 \cite{aliyari2021improving}; 181,32; 182,16 \cite{chen2019characterization}; \newline 202,32 \cite{mohammed2024improving}; 193,8 \cite{zeleke2024impact} & 193,80 \\
        \bottomrule
    \end{tabularx}
\end{table}

Os valores de referência foram compilados de estudos de caracterização físico-química desses óleos, tendo como apoio os parâmetros estabelecidos pela \textit{Food and Agriculture Organization} (FAO) \cite{fao_codex}.

\subsubsection{Cálculo da quantidade de NaOH}

O índice de saponificação é definido em relação ao KOH, mas os experimentos desse trabalho devem utilizar hidróxido de sódio (NaOH) para formar sabão em barra, é necessário converter esse índice para a base correspondente. A conversão é feita pela razão entre as massas molares das duas bases, conforme a Equação~\ref{eq:is_naoh}:

\begin{equation}
    IS_{NaOH} = IS \times \frac{MM_{NaOH}}{MM_{KOH}}
    \label{eq:is_naoh}
\end{equation}

\noindent em que $IS$ é o índice de saponificação em relação ao KOH (mg KOH/g de óleo), $IS_{NaOH}$ é o índice em relação ao NaOH (mg NaOH/g de óleo), $MM_{KOH} = 56,1$ e $MM_{NaOH} = 39,997$ correspondem às massas molares do KOH e do NaOH, respectivamente. Aplicando a conversão aos valores médios da literatura, obtêm-se os índices apresentados na Tabela~\ref{tab:is_naoh}.

\begin{table}[H]
    \centering
    \footnotesize
    \caption{Índices de saponificação convertidos para a base NaOH.}
    \label{tab:is_naoh}
    \begin{tabular}{lcc}
        \toprule
        \textbf{Tipo de óleo} & \textbf{Média IS de KOH} & \textbf{Média IS de NaOH} \\
        \midrule
        Óleo de coco & 255,84 & 182,40 \\
        Óleo de soja & 193,80 & 138,17 \\
        \bottomrule
    \end{tabular}
\end{table}

Para cada grama de óleo deve-se empregar a quantidade equivalente de NaOH, em miligramas, indicada na tabela anterior. Considerando que o objetivo do projeto envolve a otimização da produção, assume-se a utilização de 120~g de óleo. A massa total de hidróxido de sódio necessária é então calculada pela Equação~\ref{eq:massa_naoh}, que soma a contribuição de cada óleo da mistura:

\begin{equation}
    m_{NaOH} = \sum_{i} m_{\text{óleo},\,i} \times IS_{NaOH,\,i}
    \label{eq:massa_naoh}
\end{equation}

\noindent em que $m_{\text{óleo},\,i}$ é a massa do óleo $i$ e $IS_{NaOH,\,i}$ é o respectivo índice de saponificação na base NaOH. A Tabela~\ref{tab:naoh_proporcao} apresenta a quantidade de hidróxido de sódio necessária para saponificar 120~g de mistura, em função do percentual de óleo de coco empregado. Além da massa de NaOH, essa tabela também indica o volume total de cada amostra de 120~g, obtido a partir das densidades dos óleos utilizados, apresentadas na Tabela~\ref{tab:densidades}.

\begin{table}[H]
    \centering
    \footnotesize
    \caption{Densidade dos óleos utilizados.}
    \label{tab:densidades}
    \begin{tabular}{lc}
        \toprule
        \textbf{Óleo} & \textbf{Densidade (g/ml)} \\
        \midrule
        Óleo de coco & 0,912 \\
        Óleo de soja & 0,923 \\
        \bottomrule
    \end{tabular}
\end{table}

\begin{table}[H]
    \centering
    \footnotesize
    \caption{Quantidade de NaOH necessária para saponificar 120~g de mistura em função do percentual de óleo de coco.}
    \label{tab:naoh_proporcao}
    \begin{tabular}{ccc}
        \toprule
        \textbf{Ol. coco (\%)} & \textbf{Volume total (ml)} & \textbf{NaOH p/ sap. (g)} \\
        \midrule
        0\%   & 110,76 & 16,58 \\
        10\%  & 110,63 & 17,11 \\
        20\%  & 110,50 & 17,64 \\
        30\%  & 110,36 & 18,17 \\
        40\%  & 110,23 & 18,70 \\
        50\%  & 110,1  & 19,23 \\
        60\%  & 109,97 & 19,76 \\
        70\%  & 109,84 & 20,30 \\
        80\%  & 109,70 & 20,83 \\
        90\%  & 109,57 & 21,36 \\
        100\% & 109,44 & 21,89 \\
        \bottomrule
    \end{tabular}
\end{table}

\subsubsection{Otimização do processo}

A otimização do processo foi conduzida variando-se um fator por vez e mantendo-se os demais constantes, de modo a isolar o efeito de cada parâmetro sobre o produto final. Foram considerados como variáveis independentes a proporção entre os óleos de babaçu e de soja, a quantidade de ácido esteárico e a quantidade de etanol. Os demais parâmetros foram mantidos fixos ao longo de todas as práticas, conforme resumido na Tabela~\ref{tab:otimizacao}.

\begin{table}[H]
    \centering
    \footnotesize
    \caption{Variáveis e constantes adotadas no processo de otimização.}
    \label{tab:otimizacao}
    \begin{tabularx}{\textwidth}{l l >{\raggedright\arraybackslash}X}
        \toprule
        \textbf{Parâmetro} & \textbf{Tipo} & \textbf{Valor ou faixa} \\
        \midrule
        Proporção babaçu/soja          & Variável  & cerca de 18\% a 66\% de babaçu \\
        Quantidade de ácido esteárico  & Variável  & 1,7\% a 3,3\% da massa de óleo \\
        Quantidade de etanol           & Variável  & 7\% a 10\% do volume de óleo \\
        Massa total de óleo            & Constante & 120~g* \\
        Quantidade de água             & Constante & 25\% da massa de óleo \\
        Temperatura de reação          & Constante & cerca de 60~°C \\
        Agitação                       & Constante & mecânica e contínua** \\
        Tempo de cura                  & Constante & cerca de 7 dias \\
        \bottomrule
    \end{tabularx}

    {\footnotesize $^{*}$~Exceto na segunda prática, da qual 80~g de óleo foi usado. \newline \quad $^{**}$~A velocidade de agitação foi elevada conforme o aumento da viscosidade da massa.}
\end{table}

A quantidade de NaOH não foi tratada como variável independente: ela foi sempre determinada pelo índice de saponificação da mistura de óleos, acompanhando, portanto, a proporção entre o babaçu e a soja empregada em cada amostra.

O tempo de reação, por sua vez, também não foi fixado como uma variável controlada. Em vez de se estipular uma duração rígida, a etapa de aquecimento foi conduzida até que a mistura atingisse o ponto de consistência desejado — semelhante a um purê de batatas ou, idealmente, a uma vaselina densa —, validado em conjunto com a medição do pH. Dessa forma, o tempo de reação configura-se mais como um critério de parada do que como um fator de otimização.

A Tabela~\ref{tab:variacao_experimentos} resume como as variáveis foram alteradas ao longo das práticas, evidenciando a estratégia de modificar um fator por vez: a proporção entre os óleos nas duas primeiras aulas, o ácido esteárico na terceira e na quarta, e o etanol na quinta e na sexta.

\begin{table}[H]
    \centering
    \footnotesize
    \caption{Variação das variáveis ao longo dos experimentos.}
    \label{tab:variacao_experimentos}
    \begin{tabular}{lccc}
        \toprule
        \textbf{Experimento} & \textbf{Babaçu (\%)} & \makecell{\textbf{Ác. esteárico}\\\textbf{(\% da massa)}} & \makecell{\textbf{Etanol}\\\textbf{(\% do volume)}} \\
        \midrule
        Dia 1 — amostra 1 & 33 & 0   & 0  \\
        Dia 1 — amostra 2 & 65 & 0   & 0  \\
        Dia 2             & 18 & 0   & 0  \\
        Dia 3             & 33 & 3,3 & 0  \\
        Dia 4             & 33 & 1,7 & 0  \\
        Dia 5             & 37 & 3,3 & 10 \\
        Dia 6             & 33 & 3,0 & 7  \\
        \bottomrule
    \end{tabular}
\end{table}